\documentclass[11pt]{letter}

\usepackage[margin=0.75in]{geometry}
\usepackage{setspace}
\usepackage{hyperref}
\usepackage{xcolor}
\usepackage{enumitem}
\usepackage{parskip}

\hypersetup{
    colorlinks=true,
    urlcolor=black,
    linkcolor=black
}

\setlength{\parindent}{0pt}
\setstretch{1.05}
\date{}
\begin{document}

\begin{letter}{}

\opening{Dear Hiring Committee,}

I saw the job post on the Amazon careers website and thought I'd be a good fit for it.

I am a Master's student in Computer Science at the University of Cincinnati, graduating December 2026, looking for a role as a \textbf{Software Dev Engineer, Digital Twin} at \textbf{Amazon Leo}. With production experience building distributed backend services, AWS automation, and reliable deployment systems, I am excited about the opportunity to contribute to \textbf{Amazon Leo} and support the Power System team's work monitoring and predicting satellite power performance.

\textbf{Here are reasons that make me a great fit for the role:}

\begin{itemize}[leftmargin=*, itemsep=0.45em]
    \item \textbf{Production software development experience:} At Siemens Industry Software, I spent two years building and operating production services, automating infrastructure workflows with Python (Boto3) across AWS environments and cutting deployment time by 40\%.

    \item \textbf{Distributed systems and cloud services:} I migrated 40+ production environments to Kubernetes using ArgoCD, Helm, GitLab CI/CD, and HashiCorp Vault, and designed an Ansible based Blue-Green deployment framework that reduced production upgrade downtime by 50\% through automated fault recovery.

    \item \textbf{Telemetry, monitoring, and anomaly detection:} I executed 44 production environment upgrades in 4 months, monitoring live infrastructure health, diagnosing anomalies, and resolving issues across Linux, EC2, and databases, work that mirrors the real-time monitoring and fault escalation this role requires.

    \item \textbf{Python and data pipelines:} My NYC Mobility Forecasting project used PySpark and scikit-learn to build a predictive pipeline over a 2.3M-row time-series dataset, improving forecast accuracy 18.77\% to 16.52\% MAPE, experience directly applicable to predicting satellite energy state from orbital and activity data.

    \item \textbf{Testing, CI/CD, and operational excellence:} As Quality Coach and SonarQube Specialist, I increased automated code coverage from 35\% to 78\% by enforcing coding standards and CI/CD quality gates, and mentored 5 engineers through structured onboarding.
\end{itemize}

I am available for full-time work after December 2026. As a Master's student at the University of Cincinnati and someone who values professionalism and growth, I am enthusiastic about contributing to Amazon Leo's mission of delivering reliable connectivity to underserved communities.

Please find my CV attached below. I look forward to hearing from you. I would be excited to help build the ground services that keep a LEO satellite constellation powered and operating reliably.

{\setlength{\parskip}{0pt}
Regards,\\
Tanmay Bhuskute,\\
bhuskutv@mail.uc.edu $\mid$ tanmay.v.bhuskute@gmail.com $\mid$ +1 7328583027,\\
3424 Brookline Ave, Flat 19, Cincinnati, OH 45220.
}

\end{letter}

\end{document}
